% Distribución causal exacta del delta Ley 9; contenido conservado sin síntesis.
\section{Carácter regional de grado \texorpdfstring{\(5\)}{5} y mantisa reducida de acción}
\label{sec:l9s102:eta-ret}

El bloque central
\begin{equation}
 \mathcal B_c=\begin{pmatrix}7&2\\2&7\end{pmatrix}
 \label{eq:l9s102:bc-accion}
\end{equation}
posee autovalores \(9\) y \(5\). La longitud \(4\) de las órbitas de paso
\(3\), el rango \(12\) del ciclo dodecafásico y
\(104976/1944=54\) fijan el carácter
\begin{equation}
 \boxed{
 H_5(\alpha,\varphi)
 =\frac12\sqrt{\frac{1000\alpha}{\varphi}}-\alpha
 +\frac9{16}\alpha^2-\frac59\alpha^3
 +\frac7{48}\alpha^4-\frac1{54}\alpha^5.}
 \label{eq:l9s102:h5}
\end{equation}
La completación nonádica construye la coordenada analítica y el determinante
regional
\begin{equation}
 x_A=\frac{50\pi}{9}\alpha_{\mathrm{HMT}},
 \qquad
 D_A=e^{-2x_A}
 =\exp\!\left(-\frac{100\pi}{9}\alpha_{\mathrm{HMT}}\right).
 \label{eq:l9s102:xa-da}
\end{equation}
La microfibra aporta \(\rho_\pi=169\), y su serie absolutamente
convergente es
\begin{equation}
 \mathscr C_\pi(\alpha)
 =169\alpha^6+\frac{D_A\alpha^7}{1-D_A\alpha}.
 \label{eq:l9s102:cpi}
\end{equation}
La mantisa reducida posterior al retorno se construye entonces por
\begin{equation}
 \boxed{
 \eta_{\mathrm{ret}}
 =H_5-\frac{90}{\pi}\mathscr C_\pi(\alpha_{\mathrm{HMT}})
 \in\mathbb R_{>0}.}
 \label{eq:l9s102:eta-ret}
\end{equation}
Su miembro derecho contiene exclusivamente objetos generados con anterioridad
por APP--TRIT--TPK y por la clausura de \(\alpha_{\mathrm{HMT}}\).

